\title{Large Language Models Can Automatically Engineer Features for Few-Shot Tabular Learning}

\begin{abstract}
Large Language Models (LLMs) possess exceptional capabilities for solving complex and previously unseen reasoning tasks, positioning them as valuable assets in the field of tabular learning, which underpins numerous practical scenarios. In this study, we introduce an innovative in-context learning approach named \textsf{FeatLLM}, which utilizes LLMs as automated feature generators to construct an optimized input dataset for tabular prediction. The features synthesized by this method facilitate class likelihood estimation using a straightforward downstream machine learning algorithm, such as linear regression, enabling robust few-shot learning performance. Our \textsf{FeatLLM} framework relies exclusively on this basic predictor, alongside the features identified by the LLM, during inference. Unlike prevailing LLM-driven methodologies, \textsf{FeatLLM} negates the necessity of querying the LLM for individual data instances at inference. Additionally, the approach operates with minimal requirements—only API access to LLMs—and addresses restrictions associated with prompt lengths. Empirical results across a diverse array of tabular datasets confirm that \textsf{FeatLLM} produces superior rule sets, consistently yielding notable (averaging 10\%) improvements over competing solutions such as TabLLM and STUNT.
\end{abstract}

\section{Introduction}
Over recent years, Large Language Models (LLMs) have exhibited remarkable competence in generating contextually relevant outputs for new and unfamiliar tasks without the need for task-specific adaptation~\cite{creswell2022selection,imani2023mathprompter,taylor2022galactica}. Thanks to extensive training on large-scale text datasets, LLMs retain significant prior knowledge, often capable of substituting specialized domain expertise across multiple fields~\cite{Genomes2021,singhal2023large,wilcox2003role}, thereby facilitating automation in varied domains such as dialogue understanding~\cite{finch2023leveraging} and software repair~\cite{fan2023automated}. Notably, through the incorporation of task descriptions and limited examples within prompts, LLMs can grasp task context, enabling them to identify salient features and examples for the task at hand~\cite{brown2020language,zhao2023survey}. These qualities underscore the ability of LLMs to interpret data, suggesting their capacity to function as expert feature engineers.

Leveraging the knowledge and logical reasoning inherent to LLMs has been further explored within tabular learning. For example, domain facts—such as the heightened likelihood of certain illnesses among older populations—can substantially improve predictive modeling for disease outcomes. In this vein, recent research articulates task and feature definitions in natural language, formats the relevant data, and subsequently feeds this information into an LLM~\cite{dinh2022lift,wang2023anypredict}. Typically, the resulting framework involves either in-context learning~\cite{nam2023semi} or lightweight model adaptation techniques~\cite{dinh2022lift,hegselmann2023tabllm}. Empirical evidence indicates that LLM-derived prior knowledge substantially enhances model performance, especially in low-data scenarios, frequently surpassing traditional tabular learning baselines~\cite{hegselmann2023tabllm}.

Existing tabular learning techniques built on LLMs suffer from several shortcomings. For instance, end-to-end inference mandates at least one LLM pass per data instance, contributing to considerable computational demands. Additionally, many approaches require model fine-tuning to reach strong performance, which is problematic for LLMs lacking openly available training resources or API support for customization; this issue is particularly salient as leading LLMs recently unveiled restrict user interaction to inference APIs~\cite{anil2023palm,openai2023gpt4}. Moreover, the majority of methods are not robust to extended prompts~\cite{liu2023lost}; when tabular datasets contain numerous attributes, the text representation can become prohibitively lengthy, resulting in reduced feasibility or degraded performance. Together, these hurdles significantly impede the practical deployment of LLM-based tabular learning strategies.

The methodology we introduce departs from the conventional paradigm of using LLMs solely for direct prediction. Instead, we employ LLMs as agents for automated feature construction, harnessing their pre-trained knowledge in a more strategic manner. More precisely, we propose a new in-context learning paradigm, \textsf{FeatLLM}, aimed at uncovering the underlying ``criteria" that drive LLM decisions. As an illustration, when predicting the presence of a specific illness, the LLM can determine and articulate logical conditions involving data features that characterize the disease. These decision criteria—or rules—enable the generation of new features, which can substitute the original ones for task resolution. This reorientation from prediction to feature engineering enables reliance on a more lightweight model for subsequent inference, and eliminates the need for sophisticated machine learning models after feature extraction, thereby reducing inference latency. Leveraging the in-context learning strengths of LLMs, feature extraction is possible simply by appending exemplar cases to the prompt, without explicit model training. Furthermore, by utilizing ensemble strategies, such as feature bagging~\cite{breiman2001random}, the method addresses prompt length limitations effectively.

\textsf{FeatLLM} organizes the feature engineering prompt into two distinct stages: (1) analyzing the problem to discern how features relate to targets; and (2) formulating discriminative rules for class separation, informed by the insights and limited examples provided earlier. This deliberate reasoning process assists LLMs in discarding irrelevant features during rule creation. The extracted rules are implemented on the data by interpreting them as code, which results in converting each sample into a binary feature indicating compliance with individual rules. Subsequently, a simple model is fitted to predict class probabilities using these new feature representations. Robustness is further enhanced by employing ensembles, repeatedly generating features alongside feature bagging, which helps control variability. Adjusting the numbers of features and samples involved in bagging allows for effective management of prompt size constraints. \textsf{FeatLLM} eliminates the need for training and achieves efficient inference, making it suitable for deployment across diverse real-world contexts.

\textsf{FeatLLM} is tested on 13 tabular datasets with limited examples per class, demonstrating both reliability and strong performance. We provide evidence that LLMs capably combine prior knowledge with information extracted from the data to produce high-quality features. Our method consistently surpasses state-of-the-art few-shot learning approaches across numerous scenarios. The code is publicly available via an anonymized GitHub repository at~\texttt{https://github.com/Sungwon-Han/FeatLLM}.

\section{Related Work}

\subsection{Few-Shot Learning with Tabular Data} The development of few-shot learning methods has made it possible to derive generalizable features from datasets despite very limited annotated examples~\cite{chen2018closer}. Originally developed for image data, few-shot learning now exhibits substantial adaptability across various modalities, such as text, audio, and more recently, tabular data~\cite{majumder2022few,nam2023stunt,schick2021exploiting}. Relying on a small number of labeled instances, however, increases susceptibility to spurious associations~\cite{bartlett2021deep}. To address this, contemporary research integrates supplementary data sources or leverages domain-tailored prior knowledge to embed suitable inductive biases during model development. For instance, using a significant amount of unlabeled data combined with targeted augmentation for semi-supervised learning~\cite{ucar2021subtab,yoon2020vime}, or initializing models with task-agnostic unsupervised pre-training~\cite{somepalli2021saint}. Unsupervised learning objectives may include strategies such as masking or perturbing portions of the data followed by reconstruction-based training~\cite{arik2021tabnet,majmundar2022met}, or data augmentation intended for contrastive learning~\cite{bahri2022scarf,chen2020simple}. Nevertheless, due to the complex and variable characteristics of tabular data, universally effective augmentation methods remain elusive, and these tactics have shown limited success within few-shot scenarios~\cite{nam2023stunt}.

Emerging research now recommends model pre-training across a diverse collection of tabular datasets, including those unrelated to the target task, enabling subsequent transfer learning~\cite{zhang2023generative,levin2022transfer,zhu2023xtab}. For example, TransTab~\cite{wang2022transtab} leverages transformer architectures to capture semantic links among columns, utilizing column names in addition to cell values. The acquired knowledge across assorted tables and columns supports effective zero-shot or few-shot predictions in novel tabular applications. Another method, TabPFN~\cite{hollmann2022tabpfn}, synthesizes data from varied distributions for pre-training a Bayesian neural network, permitting inference after pre-training by submitting both training and test instances without further optimization. Insights accumulated from multiple datasets not only outperformed conventional tree-based techniques but were particularly advantageous in limited-label settings.

\subsection{Language-Interfaced Tabular Learning} Leveraging prior knowledge through large language models (LLMs) presents an additional promising direction~\cite{anil2023palm,openai2023gpt4}. LLMs, having been trained on vast and heterogeneous textual datasets, are capable of transferring knowledge to unencountered tasks and have proven their applicability across diverse fields~\cite{brown2020language}. Recent investigations have sought to utilize the comprehensive prior knowledge embedded in LLMs for tabular data tasks, typically by converting table data into text format and incorporating it within LLM input prompts~\cite{dinh2022lift,hegselmann2023tabllm,wang2023anypredict}. These prompts, which contain tabular content, feature explanations, and task instructions, enable the model to discern relationships between features and tasks for inference. Progress in this area has manifested as either the application of parameter-efficient training methods, including LoRA~\cite{hu2021lora} and IA3~\cite{liu2022few}, or the exploitation of in-context learning, whereby a few labeled examples are added to prompts to facilitate inference without additional model training~\cite{dinh2022lift,hegselmann2023tabllm,nam2023semi,slack2023tablet}.

Although the methods discussed above have demonstrated efficacy in enhancing few-shot tabular learning, their reliance on the LLM for every inference operation introduces practical obstacles, particularly for deployment within industry settings. The present work diverges from the traditional end-to-end, black-box framework in which each sample's prediction is fully processed by the LLM. Rather, it directs the LLM to explicitly deduce the conditions or logical rules pertaining to individual classes. With \textsf{FeatLLM}, these inferred rules are translated into programmatic code that generates new features, thus facilitating prediction without requiring ongoing LLM access. This approach exploits in-context learning, leveraging both prior knowledge and the limited information available from few-shot samples to extract relevant rules.

\section{Methods}
\textsf{FeatLLM} is engineered to derive ``rules"—that is, the specific conditions linked to each answer class—by analyzing a small collection of input samples, as illustrated in Figure~\ref{fig:main_model}. Rules are formulated through the LLM, then transformed into binary features for each sample, which indicate whether a particular rule holds for that instance. These binary features are subsequently utilized in a dot product computation with trainable, non-negative weights to predict class membership probabilities. Model training entails optimizing these weights to assess the significance of each feature in contributing to class prediction (see Section~\ref{Sec:training}). To further strengthen predictive robustness, this workflow incorporates multiple iterations using bagging, followed by ensemble integration. The subsequent sections elaborate on the problem setup and describe each procedural step in detail.

\vspace*{-0.12in}
\paragraph{Problem formulation.} Consider a tabular dataset composed of $N$ labeled instances, $\mathcal{D} = \{(\mathbf{x}^i, \mathbf{y}^i)\}_{i=1}^{N}$. Here, each data sample $\mathbf{x}^i$ is a $d$-dimensional input vector corresponding to $d$ features; the dataset $\mathcal{D}$ is accompanied by the set of natural language feature names $F = \{f_j\}_{j=1}^{d}$, with additional definitions for each feature supplied separately. In a classification context, the label $\mathbf{y}^i$ takes a value from a predefined set of categories. For $k$-shot learning scenarios, only $k$ samples ($k < N$) are randomly selected as the model's training data.

\subsection{Prompt Design for Extracting Rules}
\label{Sec:prompt_design}
To facilitate rule extraction in a manner that encourages accurate logical reasoning, we steer the language model’s approach to emulate expert handling of tabular learning tasks. The prompt construction consists of three principal elements (see Fig.~\ref{fig:prompt_for_rule} and Appendix~\ref{sec:example_rule}):

\vspace*{-0.12in}
\paragraph{Basic information description.} This segment offers fundamental information necessary for problem solving (see the orange-highlighted section in Fig.~\ref{fig:prompt_for_rule}). It encompasses: (1) a description of the task to be accomplished, specifying label information; (2) descriptions for each feature, including their types and definitions where relevant; and (3) illustrative examples. Task descriptions are phrased as questions (for example, ``Does this patient's myocardial infarction complications data indicate chronic heart failure? Yes or no?'', with further sample formulations listed in Appendix~\ref{sec:task_desc}). Feature descriptions clarify the value type and, when pertinent, provide a definition; categorical features may include enumeration of possible categories. Illustrative demonstrations serialize a subset of training examples, pairing them with their correct labels. Serialization adopts the textual format utilized in prior works~\cite{dinh2022lift,hegselmann2023tabllm}:

\begin{align}
&\text{Serialize}(\mathbf{x}^i,\mathbf{y}^i, F) = \nonumber \\ 
&``\ f_1\ \text{is}\ \mathbf{x}^i_1.\ ... \ f_d\  \text{is}\  \mathbf{x}^i_d.\ \ \text{Answer:}\  \mathbf{y}^i\ ”
\end{align}

\vspace*{-0.12in}
\paragraph{Reasoning instruction.} Our objective is to strengthen the reasoning capabilities of the LLM by introducing explicit reasoning instructions (refer to blue highlighted sections in Fig.~\ref{fig:prompt_for_rule}). The reasoning instruction opens with a statement patterned after the chain-of-thought paradigm~\cite{wei2022chain}. Subsequently, the inference process for the LLM is bifurcated into two stages. Initially, the LLM is prompted to infer causality or associations between relevant features and the provided task description, independently of any sample demonstrations; this step emphasizes reliance on background knowledge or intuitive judgment, thereby enabling the model to leverage its inherent understanding. In the next stage, the LLM integrates both exemplar demonstrations and insights gained in the prior step to formulate a set of classification rules for each target class. This layered approach helps minimize the likelihood of the model picking up false associations with irrelevant columns and supports concentrating on impactful features. To maintain a balance between specificity and prompt length, and avoid issues such as underfitting or prompt bloating, the number of rules extracted per class is capped at a fixed value (namely, 10)\footnote{Details regarding rule quantity are discussed in Section~\ref{sec:analysis}.}. Furthermore, in crafting these rules, the LLM consults the feature type information presented in the introductory description to avert type inconsistencies.

\vspace*{-0.12in}
\paragraph{Response instruction.} For improved interpretation and downstream processing of inferred rules, we delineate clear instructions regarding the expected format of LLM outputs (see yellow highlighted text in Fig.~\ref{fig:prompt_for_rule}). The prompt specifies precise guidelines concerning the response template for each answer class (i.e., response formatting) as well as the structural requirements for individual rules (i.e., rule formatting). Such direction enables the LLM to choose an appropriate rule representation in accordance with feature types. In the absence of extra prompts, the model is permitted to link multiple rules using logical connectives such as AND or OR. An illustrative example of the final rule set is provided in Appendix~\ref{sec:outcome-rule}.

\subsection{Inferring Class Likelihood via Rules}
\label{Sec:training}

\paragraph{Feature extraction via rule parsing.} The rules derived during the previous step are used to construct novel binary features. Each feature corresponds to a particular class and indicates whether a sample adheres to that class's specified rules, taking a value of '0' or '1'. These features assist in assessing the probability that a sample is assigned to a given class. Yet, because LLM-produced rules are articulated in natural language, an effective parsing method is essential to facilitate automated feature creation. Although we require consistent response and rule formatting at generation time to simplify parsing, LLM outputs can exhibit syntactic variability and noise~\cite{kovavc2023large}. For example, the LLM may replace symbols such as ``$>$” or ``$<$” with their written equivalents, ``greater than" or ``less than," adding complexity to the parsing process.

To mitigate the difficulties posed by noisy natural language parsing, rather than engineering intricate parser code, we employ the LLM directly for this task. The LLM’s proficiency in semantic interpretation, despite syntactic deviations, allows it to synthesize concise code given clear instructions (thanks to its training on code datasets). To exploit these capabilities, we incorporate the function name, details about input and output, and the relevant parsed rules into the prompt before submitting it to the LLM. Figures~\ref{fig:prompt_for_parsing} and ~\ref{sec:example_parsing}-\ref{sec:example_parse_outcome} in the Appendix provide illustrative prompts and the outputs generated. The resulting code is executed using Python’s \texttt{exec()} utility, effecting the required data transformation.

\vspace*{-0.12in}
\paragraph{Class probability estimation.} With binary rule-derived features assigned to each class, a straightforward approach for evaluating how likely a sample belongs to a class is by tallying the number of satisfied class-specific rules (i.e., summing the respective binary features per class). Nevertheless, since individual rules and features vary in their predictive value, it is necessary to ascertain their relative importance from training data. To accomplish this, we employ a standard machine learning (ML) algorithm for each class’s binary features to determine their weights. For instance, we may apply a linear model (excluding bias). Denote the binary feature vector produced from sample $\mathbf{x}^i$ for class $k$ by $\mathbf{z}_k^i \in \{0, 1\}^R$, where $R$ refers to the count of rules for class $k$ and $c$ is the total number of classes. The class probability vector $\mathbf{p}^i$ for sample $\mathbf{x}^i$ can be determined as:

\begin{align}
\text{logit}_k^i &= \max(\mathbf{w}_k, 0) \cdot \mathbf{z}_k^i, \nonumber \\
\mathbf{p}^i &= \text{Softmax}({[\text{logit}_{1}^i, ..., \text{logit}_{c}^i]}). \label{eq:infer_prob}
\end{align}

In this context, $\mathbf{w}_k$ denotes the adjustable parameters in the linear framework, encapsulating the relevance of each individual feature. By constraining these weights within the positive domain, we prevent features produced by the LLM from diminishing the probability assigned to any class during the model's training. The logits are subsequently transformed to class probabilities utilizing the softmax function. Weight optimization is conducted via minimizing the cross-entropy loss, leveraging a limited set of training examples.

\vspace*{-0.12in}
\paragraph{Ensembling with bagging.} In the final step, the procedure is reiterated multiple times to yield a suite of inference models, whose collective output forms the ensemble prediction. For $T$ iterations, the ensemble utilizes the assembled weights $\{\mathbf{w}[t]\}_{t=1}^T$, with the predicted probabilities for each class $\mathbf{p}^i$ averaged across all models, as described in Eq.~\ref{eq:infer_prob}.

Randomness is injected into the rule derivation process for the ensemble by adopting a high temperature during LLM inference. Additionally, the ordering of the few-shot example demonstrations is randomized, acknowledging that permutation impacts LLM response, as suggested by prior research~\cite{lu2022fantastically}\footnote{Appendix~\ref{sec:diversity} provides further analysis regarding rule diversity.}. Bagging is employed to exploit variability in both feature and instance selection for each iteration, a strategy that becomes increasingly beneficial as the number of samples or features grows. This ensemble method introduces two primary strengths. Firstly, the framework gains resilience against LLM-generated incorrect rules in certain runs, as the correct rules from other iterations mitigate errors, bolstering robustness. This mechanism aligns with the self-consistency property of LLMs~\cite{wang2022self}, given that frequently inferred rules across trials tend to be reliable. Secondly, the ensemble solution mitigates constraints related to LLM prompt size; for sizable tabular datasets that exceed prompt length, bagging streamlines the input while safeguarding prediction integrity. If any single set of rules fails to encompass all feature interactions, the aggregate of weak learners across trials ensures that the limitations arising from partial feature or sample coverage are offset, sustaining comprehensive performance.

\section{Experiments}
We assess the performance of \textsf{FeatLLM} across several tabular datasets, undertaking an array of analyses to better understand its operational characteristics. Due to limitations in manuscript length, supplementary experiments—including different LLM variants, qualitative explorations, and further analyses—are provided in the Appendix.

\subsection{Performance Evaluation}
\paragraph{Datasets.}
Our empirical study employs 13 datasets for tasks in binary and multi-class classification: (1) Adult~\cite{asuncion2007uci} focuses on predicting whether an individual’s income surpasses \$50,000 per year; (2) Bank~\cite{moro2014data} predicts customer subscription to term deposits; (3) Blood~\cite{yeh2009knowledge} assesses the likelihood of blood donors returning; (4) Car~\cite{kadra2021well} predicts car quality categories; (5) Communities~\cite{misc_communities_and_crime_183} estimates crime rates in specified localities; (6) Credit-g~\cite{kadra2021well} determines whether an applicant is a good or bad credit risk; (7) Diabetes\footnote{\texttt{kaggle.com/datasets/uciml/pima-indians-diabetes-database}} predicts diabetes status; (8) Heart\footnote{\texttt{kaggle.com/datasets/fedesoriano/heart-failure-prediction}} anticipates coronary artery disease occurrence; and (9) Myocardial~\cite{misc_myocardial_infarction_complications_579} predicts the presence of chronic heart failure.

Beyond these, we incorporated recent datasets and synthetic datasets that are rarely used and were not present during LLM pre-training: (10) Cultivars forecasts grain yield in soybean cultivars\footnote{\texttt{archive.ics.uci.edu/dataset/913}}; (11) NHANES classifies individuals by age group\footnote{\texttt{archive.ics.uci.edu/dataset/887}}; (12) Sequence-type assigns sequences to one of four categories—arithmetic, geometric, Fibonacci, or Collatz; (13) Solution-mix predicts whether the concentration percentage in mixtures of four solutions exceeds 0.5. The first two datasets were recently added to the UCI Machine Learning Repository, post-September 2023, guaranteeing their absence from the gpt-3.5-turbo-0613 LLM API pre-training. The last two are constructed synthetically, precluding memorization by the LLM API, allowing unbiased evaluation.

Details on dataset size and complexity, along with descriptive attributes, are summarized in Table~\ref{Tab:basic-desc}.
Each dataset features unambiguous attribute names and definitions. Some datasets, such as `Communities' and `Myocardial', contain over 100 features, which introduces additional challenges related to LLMs’ context window limitations.

\vspace*{-0.12in}
\paragraph{Baselines.}
We benchmark against nine baseline methods. The first trio comprises standard tabular classification models: (1) Logistic regression (LogReg), (2) XGBoost~\cite{chen2016xgboost}, and (3) RandomForest~\cite{ho1995random}. The following two approaches (4) SCARF~\cite{bahri2022scarf} and (5) STUNT~\cite{nam2023stunt} assume access to unlabeled data, though acquiring substantial volumes of such data is often impractical in real-world settings. The sixth baseline, (6) TabPFN~\cite{hollmann2022tabpfn}, leverages synthetic datasets with varied distributions for model pretraining. The final group comprises LLM-utilizing methods: (7) In-context learning (In-context)~\cite{wei2022emergent}, (8) TABLET~\cite{slack2023tablet}, and (9) TabLLM~\cite{hegselmann2023tabllm}. In-context learning inserts few-shot examples into prompts without additional tuning. TABLET enriches prompt construction with external classifier-generated rule-sets and prototypes to improve inference. TabLLM adopts parameter-efficient adaptation, notably IA3~\cite{liu2022few}, for fine-tuning the LLM with a limited number of labeled samples.

\vspace*{-0.12in}
\paragraph{Implementation details.}
\textsf{FeatLLM} adopts GPT-3.5 as the underlying LLM, but it is structured to be compatible with various LLM architectures (results for alternative backbones can be found in Appendix~\ref{sec:backbones}). LLM inference is performed using a temperature of 0.5 and a top-p setting of 1, which aligns with the API's default. For feature extraction, 20 ensembles are deployed and the extraction process is configured to produce 10 rules. The effects of hyper-parameter selection are depicted in Figure~\ref{fig:hyperparameter2} and discussed further in Appendix~\ref{sec:hyper-parameter}. For the linear model, we utilize the Adam optimizer with a learning rate set at 0.01, trained over 200 epochs. To identify the optimal number of epochs, we implement $k$-fold cross-validation. In comparison baselines, GPT-3.5 is leveraged for in-context learning scenarios, and T0~\cite{sanh2022multitask} is used within TabLLM according to its original experimental setup. TabLLM is constrained to LLMs with accessible public checkpoints, thereby excluding GPT-3.5 from its use. For classical machine learning models—including LogReg, XGBoost, and RandomForest—parameter selection was achieved through a combination of Grid-search and $k$-fold cross-validation. Here, the value for $k$ is assigned as either 2 or 4, ensuring each class is represented within the training subset. Further implementation details are provided in Appendix~\ref{sec:baselines}.

\vspace*{-0.12in}
\paragraph{Results.}
Tables~\ref{tab:main1} and \ref{tab:main2} display comparative results on several datasets. For each experimental setting, we conduct three independent runs, each with a distinct split for training and test sets, reporting the mean and standard deviation of AUC scores. In most cases, \textsf{FeatLLM} either leads or holds the second highest rank relative to other baseline approaches. Remarkably, our method matches the results of standard tabular baselines requiring 32 shots, with just 4 training examples (refer to Fig.~\ref{fig:compares}). Although STUNT and TabPFN achieve strong results on certain datasets, their performance is not consistently superior to classical machine learning baselines overall. In addition, approaches based on LLMs such as in-context learning, TABLET, and TabLLM, encounter difficulties when applied to datasets with a large feature space. In contrast, our method, leveraging both bagging and ensemble strategies, maintains robust performance even for datasets with many features. Notably, the bagging and ensemble methodology we employ could be ported to other LLM-based techniques; however, this would necessitate doubling inference passes for each instance, substantially increasing computational burden and thus rendering such strategies inefficient in practice.

\subsection{Analysis \& Discussion}
\label{sec:analysis}
\paragraph{Ablation study.} To elucidate the contribution of core framework components, we conduct ablation experiments targeting the following aspects: (1) \textsf{-Tuning}: removing the parameter tuning in the linear classifier, instead simply aggregating the binary features for the logit; (2) \textsf{-Ensemble}: omitting the ensemble step; (3) \textsf{-Description}: excluding the use of feature descriptions; and (4) \textsf{-Reasoning}: eliminating the Step-1 process from the reasoning instruction during rule generation.

Table~\ref{tab:ablation} reports the average change in AUC for each ablation, aggregated across datasets. Disabling any component leads to a noticeable decrease in model performance. Notably, as the number of training shots increases, weight tuning delivers greater benefits, indicating that more training data enable more precise estimation of rule importances, enhancing tuning effectiveness. Conversely, feature descriptions and reasoning instructions have their most pronounced impact when the number of shots is limited, suggesting that capitalizing on LLM prior knowledge is especially advantageous in low-data regimes. Importantly, integrating the ensemble approach yields consistent performance enhancements across all experimental configurations.

\vspace*{-0.12in}
\paragraph{Training \& inference time comparison.} We evaluate the efficiency of various approaches with respect to their training and inference times. Our experiments use the Adult dataset with a 16-shot training set configuration. For computational resources, a single A100 GPU is employed by default; however, TabLLM is run on four A100 GPUs to enable model parallelism. Table~\ref{tab:cost} details these timing metrics. One key observation is that \textsf{FeatLLM} delivers inference times that are on par with traditional approaches. Conversely, LLM-driven techniques like In-context learning, TABLET, and TabLLM incur significantly longer inference times. When comparing training durations, \textsf{FeatLLM} aligns closely with other LLM-based frameworks, only needing 30 API queries; this minimal requirement neither strains financial resources nor hinders parallel execution.

\vspace*{-0.12in}
\paragraph{Quality of features generated by \textsf{FeatLLM}.}
To assess the informativeness of the rule-based features produced by \textsf{FeatLLM} for real-world tasks, we carry out a straightforward experiment. For each dataset, we first calculate the mean AUROC between the newly constructed features and the target labels, and then identify the proportion of features whose AUROC exceeds 0.5. The outcomes are presented in Table~\ref{tab:quality}. The findings illustrate that most of the generated rules are pertinent to the target variable and offer valuable input for tackling the designated task (see ``Before tuning"). Furthermore, once the linear model undergoes tuning on the data, rules deemed less relevant (i.e., rules whose learned weights fall below a set threshold) are systematically excluded (see ``After tuning"). The threshold, set at 0.1, reflects the underlying distribution of feature weights as revealed by the histogram.

\vspace*{-0.12in}
\paragraph{Effect of adding spurious correlations.} In order to evaluate how well different approaches mitigate the influence of extraneous spurious correlations in a low-shot context, we carried out a targeted assessment. Specifically, we worked with the Heart dataset and, for the experiment, supplemented it with randomly chosen columns from the Adult dataset, which have no relevance to the Heart dataset's classification task. The models' performance was tracked as the number of non-informative columns was incrementally increased. To ensure an equitable assessment across methods, we uniformly allocated 8 labeled samples per model and omitted the use of an unlabeled dataset, thus not including approaches such as SCARF or STUNT in this evaluation.

Performance variations attributable to spurious correlations are depicted in Figure~\ref{fig:spurious}. Among the methods considered, \textsf{FeatLLM} experienced the smallest decline in performance as noise increased. This can be attributed to the framework’s emphasis on integrating prior knowledge to link features directly to the task when generating rules, which discourages reliance on irrelevant columns and, in turn, bolsters robustness. In fact, less than 2\% of the extracted rules originated from noisy columns. Nonetheless, we also noticed that \textsf{FeatLLM} may inadvertently pick up on spurious correlations and confound task-relevant causal associations when extraneous columns are sourced from datasets with some genuine connection to the Heart dataset (refer to Appendix~\ref{sec:spurious-appendix}).

\vspace*{-0.12in}
\paragraph{Hyper-parameter analysis}
The two principal hyper-parameters for \textsf{FeatLLM} are the ensemble count and the number of rules per class obtained from each LLM query. To evaluate their influence on overall performance, we ran experiments varying these parameters. We observed that increasing the ensemble count leads to progressive performance enhancement, but these improvements level off after approximately 20 ensembles (refer to Fig.~\ref{fig:appendix-hyperparameter1} in Appendix). Regarding rule quantity, statistical analysis did not reveal substantial differences. Nonetheless, setting the number of extracted rules at 10 provided the best balance between prompt length and predictive accuracy (see Fig.~\ref{fig:hyperparameter2}). We hypothesize that generating more rules increases the input's dimensionality, which can both boost available information and potentially introduce overfitting when the training data is scarce.

\section{Conclusion}
We introduce \textsf{FeatLLM}, which establishes a new benchmark for few-shot tabular learning with LLMs. Unlike approaches relying solely on LLM-driven prediction at the sample level, \textsf{FeatLLM} leverages LLMs to construct predictive criteria akin to feature engineering. By generating rules and utilizing a bagged ensemble strategy, \textsf{FeatLLM} maintains minimal inference demands, operates purely with API access (without explicit model training), and is not constrained by the feature dimensionality of the data. This approach yields strong predictive results and represents a robust, efficient alternative for deploying LLMs on practical tabular datasets.

\textsf{FeatLLM} is specifically tailored for the low-shot learning context. Moving forward, we aim to enhance its feature crafting abilities for datasets with greater sample sizes, diversify the types of features generated beyond rules, and further develop its interpretability—enabling broad applicability across numerous domains.

\section{Broader Impact}
\textsf{FeatLLM} seeks to harness the pre-existing knowledge encoded within LLMs to elevate performance on tabular learning tasks above standard supervised methods. By boosting data efficiency, our method is a meaningful step towards alleviating overfitting challenges caused by limited training data, incorporating supplementary prior knowledge. This framework proves particularly advantageous for practitioners in sectors like finance and healthcare, where labeling incurs significant costs or difficulties, and pretrained LLMs potentially contain domain-relevant expertise. To promote responsible adoption, an important avenue for future work is assessing how societal biases and misinformation within the LLMs’ prior knowledge may influence predictive outcomes—especially when this knowledge outweighs observed or collected labels.

\end{document}